\documentclass[10pt,a4paper]{amsart}
\usepackage[margin=19mm]{geometry}
\usepackage{amsmath,amssymb,mathtools}
\usepackage[scr=rsfs,cal=euler]{mathalfa}
\usepackage{xfrac}
\usepackage[unicode,hidelinks,bookmarks=false]{hyperref}
\newcommand{\ie}{i.e.\ }
\newcommand{\cf}{cf.\ }
\newcommand{\resp}{respectively\ }
\newcommand{\Subsection}{\subsection}
\providecommand{\nobreakdash}{\nobreak\hbox{-}}
\setlength{\emergencystretch}{2em}
\multlinegap0pt
\usepackage{enumerate}
\usepackage{amssymb}
\usepackage[shortlabels]{enumitem}
\usepackage{xcolor}
\newtheorem{theorem}{Theorem}
\newtheorem{proposition}{Proposition}
\usepackage{cleveref}
\crefname{theorem}{Theorem}{Theorems}
\crefname{proposition}{Proposition}{Propositions}
\title{Monochromatic Paths and a Topological Approach to Norine's Conjecture}
\author{Adam D\v{z}avoronok}
\date{Source: arXiv:2606.04181v3 (CC BY 4.0)}
\begin{document}
\maketitle

\noindent\emph{Source and licence.} Monochromatic Paths and a Topological Approach to Norine's Conjecture. Version arXiv:2606.04181v3 (CC BY 4.0), distributed by arXiv under the Creative Commons Attribution 4.0 International licence, http://creativecommons.org/licenses/by/4.0/, as recorded in the arXiv licence metadata for this exact version and shown on the abstract page https://arxiv.org/abs/2606.04181v3. Full licence notice: \texttt{licenses/arxiv-2606.04181-CC-BY-4.0.txt}.\par\smallskip

\noindent\textit{Editorial scope.} Counted unit: the article's proposition prop:existence-equivariant-map (source0.tex lines 201--203). The article's complete original proof of it is delivered with it (source0.tex lines 204--230, one \texttt{proof} environment of 27 lines). No mathematics is written, repaired or re-derived: the statement and the proof are the article's own lines, in the article's own notation, and no retained line is substituted or re-flowed.  Retained uncounted as the source-local context the proof needs: lines 163--165.  Nothing was omitted from any retained span, so the package carries no omission record.  No retained line contains a citation, so the wrapper carries no bibliography and no bibliography entry.  No retained line contains a figure environment, an image-inclusion command or drawing code, so no image is delivered and the package declares no asset dependency.  Editorial formatting only, in addition: an AMS \texttt{amsart} preamble that restates the article's own declarations and macros used by the retained lines, namely the environments \texttt{theorem}, \texttt{proposition} on separate counters, together with the standard packages those lines need.  The retained environments print in this document as Theorem 1, Proposition 1 in order of appearance, whereas the article prints the same items with the numbering of its own full document.  The complete native source of the article as distributed in the arXiv e-print for this exact version is bundled unchanged as \texttt{sources/202030/source0.tex}.
\begin{theorem}\label{thm:triangulated_squares}
 Let \(\Delta\) be a centrally symmetric simplicial complex such that \(|\Delta|\) is simply connected. Then every antipodal \(2\)-edge-colouring of the edges of the \(1\)-skeleton of \(\Delta\) contains a monochromatic path connecting some antipodal pair of vertices.
\end{theorem}

\begin{proposition}\label{prop:existence-equivariant-map}
Let $\Delta$ be as described in \cref{thm:triangulated_squares}, and let $\Delta_2$ be its $2$-skeleton. If there is no monochromatic path connecting a vertex to its antipode, then there exists a $\mathbb Z_2$-equivariant continuous map $g\colon |\Delta_2| \to S^1$, where $S^1$ is equipped with the antipodal involution.
\end{proposition}

\begin{proof}
First, note that $|\Delta|$ is simply connected if and only if $|\Delta_2|$ is as well. Let us denote the red components (including singleton vertices) by $R_1,R_2,\dots,R_m$ and the blue components canonically by $B_1,B_2,\dots,B_m$, so that $\tau(R_i)=B_i$. Since each vertex belongs to exactly one red and one blue component, we assign to every vertex $v$ a pair $(i(v),j(v))$ such that $v\in R_{i(v)}\cap B_{j(v)}$. Denote this function by $\kappa:V(\Delta_2)\rightarrow [m]^2$.

We now construct a vertex map $F\colon V(\Delta_2) \to \mathbb{R}^2\setminus\{0\}$. Choose $m$ distinct points $p_1,\dots,p_m$ on the upper open semicircle of the unit circle in $\mathbb{R}^2$ for instance, $p_k=(\cos\theta_k,\sin\theta_k)$ with $0<\theta_k<\pi$. Define
\[
F(v)=p_{i(v)}-p_{j(v)}.
\]
Because $i(v)\neq j(v)$ for all vertices, we have $F(v)\neq 0$.

Next, we show that $F$ extends linearly over each simplex without hitting the origin. Let $\sigma=\{v_1,v_2,v_3\}$ be a $2$-simplex. Clearly, either all three vertices lie in the same red component or all three lie in the same blue component. The two cases are symmetric, so assume that all three vertices lie in the same red component $R_c$. Then $i(v_1)=i(v_2)=i(v_3)=c$, and the image points are
\[
F(v_k)=p_c-p_{j(v_k)} \qquad (k\in\{1,2,3\}).
\]
Taking the dot product with $p_c$ gives
\[
(p_c-p_{j(v_k)})\cdot p_c = 1-p_{j(v_k)}\cdot p_c>0,
\]
because the points $p_i$ are distinct unit vectors on the semicircle. Thus all three image points lie in the open half-plane $\{x\in\mathbb R^2 : x\cdot p_c>0\}$, so their convex hull avoids the origin. Consequently, $F$ extends affinely over every simplex of $\Delta_2$ to a continuous map
\[
f\colon |\Delta_2|\to\mathbb R^2\setminus\{0\}.
\]
Moreover, this affine extension is equivariant: if \(x=\sum_r\lambda_r v_r\) in a simplex, then \(\tau(x)=\sum_r\lambda_r\tau(v_r)\) and hence \(f(\tau(x))=-f(x)\). Finally, normalising gives a continuous map
\[
g(x)=\frac{f(x)}{\|f(x)\|}\colon |\Delta_2|\to S^1.
\]
It follows that $g(\tau(x))=-g(x)$ for all $x\in|\Delta_2|$. Hence $g$ is $\mathbb Z_2$-equivariant.
\end{proof}

\end{document}
